\documentclass[12pt,a4paper]{article}

\usepackage[spanish,es-nodecimaldot]{babel}
\usepackage[utf8]{inputenc}
\usepackage[T1]{fontenc}
\usepackage{lmodern}
\usepackage{geometry}
\usepackage{microtype}
\usepackage{setspace}
\usepackage{parskip}
\usepackage{enumitem}
\usepackage{xcolor}
\usepackage{listings}
\usepackage{hyperref}
\usepackage{amsmath}
\usepackage{amssymb}
\usepackage{array}
\usepackage{booktabs}
\usepackage{longtable}

\geometry{margin=2.5cm}
\onehalfspacing
\setlist[itemize]{leftmargin=*,itemsep=0.2em}
\setlist[enumerate]{leftmargin=*,itemsep=0.2em}

\hypersetup{
  colorlinks=true,
  linkcolor=black,
  urlcolor=blue,
  pdftitle={Resolucion del Practico 5 - Testing de Software 2024},
  pdfauthor={}
}

\lstdefinestyle{codestyle}{
  basicstyle=\ttfamily\small,
  frame=single,
  breaklines=true,
  showstringspaces=false,
  keywordstyle=\color{blue!70!black}\bfseries,
  commentstyle=\color{gray!70!black}\itshape,
  stringstyle=\color{orange!70!black}
}

\title{\textbf{Resolucion del Practico 5}\\Testing de Software 2024}
\author{}
\date{}

\begin{document}

\maketitle
\tableofcontents
\newpage

\section*{Introduccion}
\addcontentsline{toc}{section}{Introduccion}

Este documento contiene la resolucion del \textit{Practico 5} de Testing de Software, centrado
en \textbf{testing basado en expresiones logicas}. Los predicados que aparecen en el codigo
(decisiones de \texttt{if}, \texttt{while}, etc.), en especificaciones, en maquinas de estados o
en requisitos, definen un espacio de combinaciones de valores de verdad mucho mas grande del que
es razonable probar. Los criterios de cobertura logica seleccionan, de manera sistematica, un
subconjunto significativo de esas combinaciones.

A lo largo de la resolucion se usan los siguientes criterios:

\begin{itemize}
  \item \textbf{PC} (Predicate Coverage): cada predicado $p$ se evalua a \textit{true} y a
  \textit{false}.
  \item \textbf{CC} (Clause Coverage): cada clausula $c$ se evalua a \textit{true} y a
  \textit{false}.
  \item \textbf{CACC} (Correlated Active Clause Coverage): para cada clausula mayor $c_i$ existe
  un par de tests donde $c_i$ \emph{determina} el predicado y, en ese par, $p$ cambia de valor.
  \item \textbf{RACC} (Restricted Active Clause Coverage): igual que CACC, pero ademas las
  clausulas menores deben tener \emph{exactamente los mismos valores} en ambos tests del par.
\end{itemize}

\textbf{Determinacion de una clausula.} Una clausula $c_i$ determina a un predicado $p$ cuando,
fijando las demas clausulas (las \emph{menores}), el valor de $p$ es el mismo que el valor de
$c_i$. Operativamente, $c_i$ determina $p$ si al cambiar $c_i$ y dejar fijas las demas, $p$
cambia.

\textbf{Convenciones.} En las tablas, las filas se numeran en orden lexicografico con $T$ antes
que $F$. Cuando se enumeran pares CACC/RACC, se listan como $(\text{fila}_i, \text{fila}_j)$.

\section{Ejercicio 1: Cobertura logica sobre cinco predicados}

Para cada uno de los predicados se identifican sus clausulas, se construye la tabla de verdad y
se analiza cada criterio: CC, PC, CACC y RACC.

\subsection{Predicado $p = a \vee (b \wedge c)$}

\paragraph{(a) Clausulas.} Las clausulas son $a$, $b$ y $c$.

\paragraph{(b) Tabla de verdad.}

\begin{center}
\begin{tabular}{ccccc}
\toprule
Fila & $a$ & $b$ & $c$ & $p$ \\
\midrule
1 & T & T & T & T \\
2 & T & T & F & T \\
3 & T & F & T & T \\
4 & T & F & F & T \\
5 & F & T & T & T \\
6 & F & T & F & F \\
7 & F & F & T & F \\
8 & F & F & F & F \\
\bottomrule
\end{tabular}
\end{center}

\paragraph{(c) CC pero no PC.} \textbf{Si es posible.} Si la suite mantiene $p$ constante pero
cada clausula cambia, se logra CC sin PC.

Suite: $\{$fila 4, fila 5$\}$ $=\{(T,F,F),\,(F,T,T)\}$. Cada clausula toma $T$ y $F$, y en ambos
casos $p = T$. Se cumple CC y no se cumple PC.

\paragraph{(d) PC pero no CC.} \textbf{Si es posible.}

Suite: $\{$fila 1, fila 6$\}$ $=\{(T,T,T),\,(F,T,F)\}$. Aqui $p$ toma $T$ y $F$ (PC), pero la
clausula $b$ se queda fija en $T$, asi que CC falla.

\paragraph{(e) Pares CACC.} Para cada clausula mayor se buscan pares en los que, fijando un valor
de las menores que permita la determinacion, el predicado cambie:

\begin{itemize}
  \item Mayor $a$: las menores deben hacer $b \wedge c = F$. Pares cruzados entre
  $\{2,3,4\}$ y $\{6,7,8\}$, es decir 9 pares en total: $(2,6), (2,7), (2,8), (3,6), (3,7),
  (3,8), (4,6), (4,7), (4,8)$.
  \item Mayor $b$: para que $b$ determine $(b\wedge c)$ debe ser $c=T$; ademas se necesita
  $a=F$ para que $a$ no fuerce $p$. Par: $(5,7)$.
  \item Mayor $c$: simetricamente, $b=T$ y $a=F$. Par: $(5,6)$.
\end{itemize}

\paragraph{(f) Pares RACC.} RACC exige que las clausulas menores tengan los \emph{mismos} valores
en ambos tests del par:

\begin{itemize}
  \item Mayor $a$: parejas con $(b,c)$ iguales. $(2,6)$ con $(b,c)=(T,F)$; $(3,7)$ con $(F,T)$;
  $(4,8)$ con $(F,F)$.
  \item Mayor $b$: $(5,7)$, con $(a,c)=(F,T)$.
  \item Mayor $c$: $(5,6)$, con $(a,b)=(F,T)$.
\end{itemize}

\subsection{Predicado $p = a \rightarrow (b \rightarrow c)$}

\paragraph{(a) Clausulas.} $a$, $b$, $c$. Equivalentemente $p \equiv \neg a \vee \neg b \vee c$.

\paragraph{(b) Tabla de verdad.}

\begin{center}
\begin{tabular}{ccccc}
\toprule
Fila & $a$ & $b$ & $c$ & $p$ \\
\midrule
1 & T & T & T & T \\
2 & T & T & F & F \\
3 & T & F & T & T \\
4 & T & F & F & T \\
5 & F & T & T & T \\
6 & F & T & F & T \\
7 & F & F & T & T \\
8 & F & F & F & T \\
\bottomrule
\end{tabular}
\end{center}

\paragraph{(c) CC pero no PC.} \textbf{Si es posible.} El predicado solo es $F$ en la fila 2; si
se evita esa fila, $p$ queda en $T$ y aun se puede tener CC.

Suite: $\{$fila 1, fila 8$\}$ $=\{(T,T,T),\,(F,F,F)\}$. Cada clausula toma $T$ y $F$, $p$ es $T$
en ambos casos. CC se cumple, PC no.

\paragraph{(d) PC pero no CC.} \textbf{Si es posible.}

Suite: $\{$fila 1, fila 2$\}$ $=\{(T,T,T),\,(T,T,F)\}$. $p$ toma ambos valores; $a$ y $b$ quedan
fijos en $T$.

\paragraph{(e) Pares CACC.}

\begin{itemize}
  \item Mayor $a$: para que $a$ determine $p$, el subpredicado $(b\rightarrow c)$ debe ser $F$,
  o sea $b=T,\, c=F$. Par: $(2,6)$.
  \item Mayor $b$: $b$ determina cuando $a=T$ (sino $p$ es $T$ siempre) y $c=F$. Par: $(2,4)$.
  \item Mayor $c$: $c$ determina cuando $a=T$ y $b=T$. Par: $(1,2)$.
\end{itemize}

\paragraph{(f) Pares RACC.} Como las menores quedan fijadas a un unico valor, los pares CACC ya
son RACC: $(2,6)$ para $a$, $(2,4)$ para $b$, $(1,2)$ para $c$.

\subsection{Predicado $p = (a \vee b) \wedge (c \vee d)$}

\paragraph{(a) Clausulas.} $a$, $b$, $c$, $d$.

\paragraph{(b) Tabla de verdad.}

\begin{center}
\begin{tabular}{cccccc}
\toprule
Fila & $a$ & $b$ & $c$ & $d$ & $p$ \\
\midrule
1  & T & T & T & T & T \\
2  & T & T & T & F & T \\
3  & T & T & F & T & T \\
4  & T & T & F & F & F \\
5  & T & F & T & T & T \\
6  & T & F & T & F & T \\
7  & T & F & F & T & T \\
8  & T & F & F & F & F \\
9  & F & T & T & T & T \\
10 & F & T & T & F & T \\
11 & F & T & F & T & T \\
12 & F & T & F & F & F \\
13 & F & F & T & T & F \\
14 & F & F & T & F & F \\
15 & F & F & F & T & F \\
16 & F & F & F & F & F \\
\bottomrule
\end{tabular}
\end{center}

\paragraph{(c) CC pero no PC.} \textbf{Si es posible.}

Suite: $\{$fila 4, fila 13$\}$ $=\{(T,T,F,F),\,(F,F,T,T)\}$. Cada clausula cambia, pero $p=F$
en ambos casos.

\paragraph{(d) PC pero no CC.} \textbf{Si es posible.}

Suite: $\{$fila 1, fila 4$\}$ $=\{(T,T,T,T),\,(T,T,F,F)\}$. $p$ toma $T$ y $F$ pero $a$ y $b$
quedan fijas en $T$.

\paragraph{(e) Pares CACC.}

\begin{itemize}
  \item Mayor $a$: las menores deben hacer $b=F$ y $(c\vee d)=T$. Filas con $a=T$ que cumplen
  esas condiciones: $5,6,7$. Filas con $a=F$ que las cumplen: $13,14,15$. Pares: cualquier
  combinacion cruzada (9 pares).
  \item Mayor $b$: simetrica a $a$ (con $a=F$, $(c\vee d)=T$): filas $9,10,11$ frente a $13,14,
  15$. 9 pares en total.
  \item Mayor $c$: necesita $d=F$ (sino $c\vee d=T$ siempre) y $(a\vee b)=T$. Filas con $c=T$ que
  cumplen $d=F$ y $(a\vee b)=T$: $2,6,10$. Filas con $c=F$ en las mismas condiciones: $4,8,12$.
  9 pares.
  \item Mayor $d$: por simetria con $c$, $c=F$ y $(a\vee b)=T$. Filas con $d=T$: $3,7,11$;
  filas con $d=F$: $4,8,12$. 9 pares.
\end{itemize}

\paragraph{(f) Pares RACC.} Se pide que las menores valgan lo mismo en cada par:

\begin{itemize}
  \item Mayor $a$ (fijando $(b,c,d)$): $(5,13)$ con $(F,T,T)$; $(6,14)$ con $(F,T,F)$; $(7,15)$
  con $(F,F,T)$.
  \item Mayor $b$ (fijando $(a,c,d)$): $(9,13)$, $(10,14)$, $(11,15)$.
  \item Mayor $c$ (fijando $(a,b,d)$): $(2,4)$, $(6,8)$, $(10,12)$.
  \item Mayor $d$ (fijando $(a,b,c)$): $(3,4)$, $(7,8)$, $(11,12)$.
\end{itemize}

\subsection{Predicado $p = a \oplus b$}

\paragraph{(a) Clausulas.} $a$, $b$.

\paragraph{(b) Tabla de verdad.}

\begin{center}
\begin{tabular}{cccc}
\toprule
Fila & $a$ & $b$ & $p$ \\
\midrule
1 & T & T & F \\
2 & T & F & T \\
3 & F & T & T \\
4 & F & F & F \\
\bottomrule
\end{tabular}
\end{center}

\paragraph{(c) CC pero no PC.} \textbf{Si es posible.}

Suite: $\{$fila 1, fila 4$\}$ $=\{(T,T),\,(F,F)\}$. $p=F$ en ambos casos, pero ambas clausulas
toman $T$ y $F$.

\paragraph{(d) PC pero no CC.} \textbf{Si es posible.}

Suite: $\{$fila 1, fila 2$\}$ $=\{(T,T),\,(T,F)\}$. $p$ toma ambos valores, pero $a$ no toma $F$.

\paragraph{(e) y (f) Pares CACC y RACC.} Para XOR, cada clausula siempre determina al predicado
(cambiar una variable siempre cambia $p$). Ademas, fijando el valor de la otra clausula los
pares son unicos, asi que CACC y RACC coinciden:

\begin{itemize}
  \item Mayor $a$: $(1,3)$ con $b=T$; $(2,4)$ con $b=F$.
  \item Mayor $b$: $(1,2)$ con $a=T$; $(3,4)$ con $a=F$.
\end{itemize}

\subsection{Predicado $p = a \vee b \vee (c \wedge d)$}

\paragraph{(a) Clausulas.} $a$, $b$, $c$, $d$.

\paragraph{(b) Tabla de verdad.}

\begin{center}
\begin{tabular}{cccccc}
\toprule
Fila & $a$ & $b$ & $c$ & $d$ & $p$ \\
\midrule
1  & T & T & T & T & T \\
2  & T & T & T & F & T \\
3  & T & T & F & T & T \\
4  & T & T & F & F & T \\
5  & T & F & T & T & T \\
6  & T & F & T & F & T \\
7  & T & F & F & T & T \\
8  & T & F & F & F & T \\
9  & F & T & T & T & T \\
10 & F & T & T & F & T \\
11 & F & T & F & T & T \\
12 & F & T & F & F & T \\
13 & F & F & T & T & T \\
14 & F & F & T & F & F \\
15 & F & F & F & T & F \\
16 & F & F & F & F & F \\
\bottomrule
\end{tabular}
\end{center}

\paragraph{(c) CC pero no PC.} \textbf{Si es posible.}

Suite: $\{$fila 4, fila 13$\}$ $=\{(T,T,F,F),\,(F,F,T,T)\}$. Cada clausula cambia y $p=T$ en
ambos casos.

\paragraph{(d) PC pero no CC.} \textbf{Si es posible.}

Suite: $\{$fila 1, fila 14$\}$ $=\{(T,T,T,T),\,(F,F,T,F)\}$. $p$ toma ambos valores; $c$ queda
fija en $T$.

\paragraph{(e) Pares CACC.}

\begin{itemize}
  \item Mayor $a$: necesita $b=F$ y $(c\wedge d)=F$. Filas con $a=T$: $6,7,8$; filas con $a=F$:
  $14,15,16$. 9 pares cruzados.
  \item Mayor $b$: simetrica a $a$ (con $a=F$ y $(c\wedge d)=F$): filas $10,11,12$ frente a
  $14,15,16$. 9 pares.
  \item Mayor $c$: $c$ determina $(c \wedge d)$ cuando $d=T$. Ademas, para que $(c\wedge d)$
  determine $p$, se necesita $a=F$ y $b=F$. Par: $(13,15)$.
  \item Mayor $d$: simetrico a $c$, con $c=T$, $a=F$, $b=F$. Par: $(13,14)$.
\end{itemize}

\paragraph{(f) Pares RACC.} Fijando las menores en cada par:

\begin{itemize}
  \item Mayor $a$ (fijando $(b,c,d)$): $(6,14)$ con $(F,T,F)$; $(7,15)$ con $(F,F,T)$; $(8,16)$
  con $(F,F,F)$.
  \item Mayor $b$ (fijando $(a,c,d)$): $(10,14)$, $(11,15)$, $(12,16)$.
  \item Mayor $c$: $(13,15)$.
  \item Mayor $d$: $(13,14)$.
\end{itemize}

\section{Ejercicio 2: \texttt{checkIt} y \texttt{checkItExpand}}

\subsection{Codigo de partida}

\begin{lstlisting}[style=codestyle,language=Java]
public static void checkIt (boolean a, boolean b, boolean c) {
    if (a && (b || c))
        System.out.println ("P is true");
    else
        System.out.println ("P isn't true");
}
\end{lstlisting}

El predicado bajo analisis es $p = a \wedge (b \vee c)$.

\subsection{(a) Transformacion a \texttt{checkItExpand} e instrumentacion}

Se reescribe el codigo para que \emph{cada \texttt{if} use una sola variable booleana}. Una forma
estandar es descomponer el \texttt{if} en \texttt{if/else} anidados:

\begin{lstlisting}[style=codestyle,language=Java]
public static ExecutionTrace checkItExpand(boolean a, boolean b, boolean c) {
    nodes.add("START"); nodes.add("A");
    if (a) {
        nodes.add("B");
        if (b) {
            nodes.add("TRUE"); result = true;
        } else {
            nodes.add("C");
            if (c) { nodes.add("TRUE"); result = true; }
            else   { nodes.add("FALSE"); result = false; }
        }
    } else {
        nodes.add("FALSE"); result = false;
    }
    nodes.add("END");
    return new ExecutionTrace(result, nodes);
}
\end{lstlisting}

\textbf{Instrumentacion.} \texttt{checkItExpand} devuelve un objeto \texttt{ExecutionTrace} que
contiene:

\begin{itemize}
  \item la lista de \emph{nodos visitados} en orden;
  \item las \emph{aristas atravesadas}, derivadas como pares consecutivos
  $N_i \to N_{i+1}$.
\end{itemize}

Los nodos del CFG son
$\{\textbf{START},\textbf{A},\textbf{B},\textbf{C},\textbf{TRUE},\textbf{FALSE},\textbf{END}\}$ y
las aristas posibles:
\[
\textbf{START}{\to}\textbf{A},\;
\textbf{A}{\to}\textbf{B},\;
\textbf{A}{\to}\textbf{FALSE},\;
\textbf{B}{\to}\textbf{TRUE},\;
\textbf{B}{\to}\textbf{C},\;
\textbf{C}{\to}\textbf{TRUE},\;
\textbf{C}{\to}\textbf{FALSE},\;
\textbf{TRUE}{\to}\textbf{END},\;
\textbf{FALSE}{\to}\textbf{END}.
\]

La instrumentacion permite verificar cobertura de aristas directamente desde los tests, sin
depender de una herramienta externa.

\subsection{(b) Suites T1 (CACC) y T2 (Edge Coverage). \textquestiondown{}T2 satisface CACC?}

\paragraph{Suite T1 (CACC sobre \texttt{checkIt}).}

\[
T_1 = \{ t_1=(T,T,F),\; t_2=(F,T,F),\; t_3=(T,F,F),\; t_4=(T,F,T) \}.
\]

Justificacion por clausula mayor:

\begin{itemize}
  \item Mayor $a$: $(t_1,t_2)=((T,T,F),(F,T,F))$. Las menores $b\vee c$ valen $T$ en ambos casos;
  $p$ cambia de $T$ a $F$.
  \item Mayor $b$: $(t_1,t_3)=((T,T,F),(T,F,F))$. $a=T$ y $c=F$ fijos; $p$ cambia de $T$ a $F$.
  \item Mayor $c$: $(t_4,t_3)=((T,F,T),(T,F,F))$. $a=T$ y $b=F$ fijos; $p$ cambia de $T$ a $F$.
\end{itemize}

\paragraph{Suite T2 (Edge Coverage sobre \texttt{checkItExpand}).}

\[
T_2 = \{ e_1=(F,F,F),\; e_2=(T,T,F),\; e_3=(T,F,T),\; e_4=(T,F,F) \}.
\]

Caminos:

\begin{itemize}
  \item $e_1$: START $\to$ A $\to$ FALSE $\to$ END.
  \item $e_2$: START $\to$ A $\to$ B $\to$ TRUE $\to$ END.
  \item $e_3$: START $\to$ A $\to$ B $\to$ C $\to$ TRUE $\to$ END.
  \item $e_4$: START $\to$ A $\to$ B $\to$ C $\to$ FALSE $\to$ END.
\end{itemize}

La union cubre las 9 aristas. \textbf{T2 satisface Edge Coverage.}

\paragraph{\textquestiondown{}T2 satisface CACC sobre \texttt{checkIt}?} \textbf{No.}

Para que $a$ sea la clausula mayor en \texttt{checkIt}, hace falta un par donde
$b \vee c = T$ y $a$ cambie entre los dos tests. En $T_2$ los casos con $b \vee c = T$ son $e_2$
y $e_3$, ambos con $a=T$. \emph{No hay ningun caso con $a=F$ y $b\vee c = T$.} Por lo tanto, el
par CACC para mayor $a$ no existe en $T_2$.

Edge Coverage sobre el CFG expandido y CACC sobre el predicado original son criterios
distintos: la expansion estructural conserva la semantica del programa pero pierde la nocion de
clausula activa, justamente porque cada \texttt{if} solo refiere a una variable.

\subsection{(c) Implementacion JUnit de T1 y T2}

\paragraph{T1 (CACC).} La suite \texttt{CheckItT1CaccTest} verifica, para cada clausula mayor,
que el par correspondiente produce un cambio en el predicado:

\begin{lstlisting}[style=codestyle,language=Java]
@Test public void caccMajorA() {
    assertTrue (CheckIt.checkIt(true,  true,  false));
    assertFalse(CheckIt.checkIt(false, true,  false));
}
@Test public void caccMajorB() {
    assertTrue (CheckIt.checkIt(true,  true,  false));
    assertFalse(CheckIt.checkIt(true,  false, false));
}
@Test public void caccMajorC() {
    assertTrue (CheckIt.checkIt(true,  false, true));
    assertFalse(CheckIt.checkIt(true,  false, false));
}
\end{lstlisting}

\paragraph{T2 (Edge Coverage).} La suite \texttt{CheckItT2EdgeCoverageTest} ejecuta
\texttt{checkItExpand} sobre los cuatro casos, junta las aristas y compara contra el conjunto
esperado:

\begin{lstlisting}[style=codestyle,language=Java]
Set<String> covered = new LinkedHashSet<>();
covered.addAll(CheckIt.checkItExpand(false, false, false).getEdges());
covered.addAll(CheckIt.checkItExpand(true,  true,  false).getEdges());
covered.addAll(CheckIt.checkItExpand(true,  false, true ).getEdges());
covered.addAll(CheckIt.checkItExpand(true,  false, false).getEdges());

assertEquals(EXPECTED_EDGES, covered);
\end{lstlisting}

Un segundo test verifica explicitamente que esa suite \emph{no} contiene ningun par CACC para
$a$, dejando documentado el resultado del inciso (b).

\paragraph{Resultado de ejecucion.} Desde \texttt{tp5/assignment-5-rodeghiero}:

\begin{lstlisting}[style=codestyle]
JAVA_HOME=$(/usr/libexec/java_home -v 17) mvn -Dmaven.repo.local=.m2 \
    -Djacoco.skip=true test
\end{lstlisting}

reporta \texttt{BUILD SUCCESS} con todos los tests pasando.

\section{Ejercicio 3: \texttt{Thermostat.turnHeaterOn}}

\subsection{Predicado bajo analisis}

La decision principal del metodo es
\[
((\textit{curTemp} < \textit{dTemp} - \textit{thresholdDiff}) \;\vee\;
  (\textit{override} \wedge \textit{curTemp} < \textit{overTemp} - \textit{thresholdDiff}))
  \;\wedge\; (\textit{timeSinceLastRun} > \textit{minLag}),
\]
que se reescribe en forma abstracta como
\[
p = ((a \vee (b \wedge c)) \wedge d),
\]
con las clausulas
\begin{itemize}
  \item $a$: \texttt{curTemp < dTemp - thresholdDiff},
  \item $b$: \texttt{override},
  \item $c$: \texttt{curTemp < overTemp - thresholdDiff},
  \item $d$: \texttt{timeSinceLastRun > minLag}.
\end{itemize}

\subsection{Diseno comun de los tests}

Para que cada test controle los valores de las cuatro clausulas, se fija:
\texttt{Period.MORNING}, \texttt{DayType.WEEKDAY},
\texttt{setSetting(MORNING,WEEKDAY,69)}, \texttt{thresholdDiff = 5} y \texttt{minLag = 10}.
Variando \texttt{curTemp}, \texttt{override}, \texttt{overTemp} y \texttt{timeSinceLastRun} se
elige el valor de cada clausula. La logica de armado esta en
\texttt{ThermostatCoverageSupport.runCase(a,b,c,d)}, que ademas devuelve los valores
\emph{reales} de cada clausula evaluados sobre el objeto para verificar que el setup coincide
con la intencion del test.

\subsection{(a) CC pero no PC}

\textbf{Suite propuesta:}
\begin{itemize}
  \item $t_1 = (a=T,b=F,c=F,d=F)$ $\Rightarrow$ $p = (T \vee F) \wedge F = F$.
  \item $t_2 = (a=F,b=T,c=F,d=T)$ $\Rightarrow$ $p = (F \vee F) \wedge T = F$.
  \item $t_3 = (a=F,b=F,c=T,d=T)$ $\Rightarrow$ $p = (F \vee F) \wedge T = F$.
\end{itemize}

\textbf{Justificacion.}

\begin{itemize}
  \item Cada clausula toma $T$ y $F$ al menos una vez:
  \begin{itemize}
    \item $a$: $T$ en $t_1$; $F$ en $t_2, t_3$.
    \item $b$: $T$ en $t_2$; $F$ en $t_1, t_3$.
    \item $c$: $T$ en $t_3$; $F$ en $t_1, t_2$.
    \item $d$: $T$ en $t_2, t_3$; $F$ en $t_1$.
  \end{itemize}
  Por lo tanto \textbf{CC se cumple}.
  \item $p$ vale $F$ en los tres casos, asi que \textbf{PC no se cumple}.
\end{itemize}

\subsection{(b) PC pero no CC}

\textbf{Suite propuesta:}
\begin{itemize}
  \item $u_1 = (a=T,b=T,c=T,d=T)$ $\Rightarrow$ $p=T$.
  \item $u_2 = (a=T,b=T,c=T,d=F)$ $\Rightarrow$ $p=F$.
\end{itemize}

$p$ toma $T$ y $F$, asi que PC se cumple. Sin embargo $a$, $b$ y $c$ no toman el valor $F$ en
ninguno de los dos tests; CC no se cumple.

\subsection{(c) CACC}

Se construye un par por cada clausula mayor.

\paragraph{Mayor $a$.} Necesita $b \wedge c = F$ (sino el subpredicado $(a \vee (b \wedge c))$
no esta determinado por $a$) y $d = T$ (sino el $\wedge$ exterior anula a $a$).
\[
(T,F,F,T) \;\Rightarrow\; p=T, \qquad (F,F,F,T) \;\Rightarrow\; p=F.
\]

\paragraph{Mayor $b$.} Para que $b$ determine a $b \wedge c$ debe ser $c=T$; ademas
$a=F$ y $d=T$.
\[
(F,T,T,T) \;\Rightarrow\; p=T, \qquad (F,F,T,T) \;\Rightarrow\; p=F.
\]

\paragraph{Mayor $c$.} Simetrica a $b$: $b=T$, $a=F$, $d=T$.
\[
(F,T,T,T) \;\Rightarrow\; p=T, \qquad (F,T,F,T) \;\Rightarrow\; p=F.
\]

\paragraph{Mayor $d$.} Necesita $(a \vee (b\wedge c)) = T$, lo cual se logra con $a=T$.
\[
(T,F,F,T) \;\Rightarrow\; p=T, \qquad (T,F,F,F) \;\Rightarrow\; p=F.
\]

Reuniendo los pares y eliminando duplicados, la suite CACC contiene 6 tests distintos.

\subsection{Implementaciones JUnit}

Los tres incisos se materializan como tres clases en
\texttt{src/test/java/assignment5\_exercises/thermostat}:
\begin{itemize}
  \item \texttt{ThermostatExercise3ClauseCoverageTest}: verifica el item (a).
  \item \texttt{ThermostatExercise3PredicateCoverageTest}: verifica el item (b).
  \item \texttt{ThermostatExercise3CaccTest}: verifica el item (c) con cuatro pares CACC.
\end{itemize}
Las tres clases comparten \texttt{ThermostatCoverageSupport}, que ademas chequea que el valor
realmente devuelto por \texttt{turnHeaterOn} coincida con la evaluacion teorica del predicado.
La suite completa pasa con \texttt{mvn test}.

\section{Ejercicio 4: \texttt{TriTyp.triang}}

\subsection{Predicados identificados}

\texttt{triang(s1,s2,s3)} contiene varios predicados que se evaluan en cascada. Los identificamos
asi:

\begin{enumerate}
  \item $P_1 = (s_1 \le 0) \vee (s_2 \le 0) \vee (s_3 \le 0)$, con clausulas $c_1, c_2, c_3$.
  \item $P_2 = (s_1 = s_2)$.
  \item $P_3 = (s_1 = s_3)$.
  \item $P_4 = (s_2 = s_3)$.
  \item $P_5 = (\textit{triOut} = 0)$.
  \item $P_6 = (s_1+s_2 \le s_3) \vee (s_2+s_3 \le s_1) \vee (s_1+s_3 \le s_2)$, clausulas
  $d_1,d_2,d_3$.
  \item $P_7 = (\textit{triOut} > 3)$.
  \item $P_8 = (\textit{triOut} = 1) \wedge (s_1+s_2 > s_3)$, clausulas $e_1, e_2$.
  \item $P_9 = (\textit{triOut} = 2) \wedge (s_1+s_3 > s_2)$, clausulas $f_1, f_2$.
  \item $P_{10} = (\textit{triOut} = 3) \wedge (s_2+s_3 > s_1)$, clausulas $g_1, g_2$.
\end{enumerate}

Para CACC distinguimos:

\begin{itemize}
  \item Predicados con varias clausulas ($P_1, P_6, P_8, P_9, P_{10}$): un par por cada clausula
  mayor.
  \item Predicados de una sola clausula ($P_2, P_3, P_4, P_5, P_7$): basta con que cada uno tome
  valor $T$ y $F$ a lo largo de la suite.
\end{itemize}

\subsection{Casos de prueba}

\begin{center}
\begin{tabular}{lll}
\toprule
$(s_1,s_2,s_3)$ & \texttt{triang} & Comentario \\
\midrule
$(-1,1,1)$ & 4 & $P_1$: $c_1$ activa \\
$(1,-1,1)$ & 4 & $P_1$: $c_2$ activa \\
$(1,1,-1)$ & 4 & $P_1$: $c_3$ activa \\
$(1,1,1)$  & 3 & $P_1$ falso, equilatero \\
$(1,2,3)$  & 4 & $P_6$: $d_1$ activa \\
$(2,3,4)$  & 1 & $P_6$ falso, escaleno valido \\
$(3,1,2)$  & 4 & $P_6$: $d_2$ activa \\
$(1,3,2)$  & 4 & $P_6$: $d_3$ activa \\
$(1,2,1)$  & 4 & isosceles imposible ($s_1+s_3=s_2$): $P_9$ falso \\
$(2,2,1)$  & 2 & isosceles via $P_8$ verdadero \\
$(1,1,2)$  & 4 & $P_8$/$P_{10}$: $e_2 = F$, $g_1=F$ \\
$(2,1,2)$  & 2 & isosceles via $P_9$ verdadero \\
$(1,2,2)$  & 2 & isosceles via $P_{10}$ verdadero \\
$(2,1,1)$  & 4 & $P_{10}$: $g_2$ falso \\
\bottomrule
\end{tabular}
\end{center}

\subsection{Pares CACC por predicado}

\paragraph{$P_1$: $(s_1\le 0)\vee(s_2\le 0)\vee(s_3\le 0)$.}
\begin{itemize}
  \item Mayor $c_1$: $(-1,1,1)$ vs $(1,1,1)$, $c_2,c_3 = F$ en ambos.
  \item Mayor $c_2$: $(1,-1,1)$ vs $(1,1,1)$, $c_1,c_3 = F$ en ambos.
  \item Mayor $c_3$: $(1,1,-1)$ vs $(1,1,1)$, $c_1,c_2 = F$ en ambos.
\end{itemize}

\paragraph{$P_6$: $(s_1+s_2 \le s_3)\vee\ldots$.}
\begin{itemize}
  \item Mayor $d_1$: $(1,2,3)$ vs $(2,3,4)$.
  \item Mayor $d_2$: $(3,1,2)$ vs $(2,3,4)$.
  \item Mayor $d_3$: $(1,3,2)$ vs $(2,3,4)$.
\end{itemize}
En cada par, las clausulas $d_i$ no mayores valen $F$ en ambos lados, asi que el cambio en
$P_6$ se debe exclusivamente a la clausula mayor.

\paragraph{$P_8$: $(\textit{triOut}=1) \wedge (s_1+s_2>s_3)$.}
\begin{itemize}
  \item Mayor $e_1$: $(1,2,1)$ vs $(2,2,1)$. En ambos $e_2=T$; $e_1$ cambia de $F$ a $T$ y $P_8$
  pasa de $F$ a $T$.
  \item Mayor $e_2$: $(2,2,1)$ vs $(1,1,2)$. En ambos $e_1=T$; $e_2$ cambia de $T$ a $F$ y $P_8$
  pasa de $T$ a $F$.
\end{itemize}

\paragraph{$P_9$: $(\textit{triOut}=2) \wedge (s_1+s_3>s_2)$.}
\begin{itemize}
  \item Mayor $f_1$: $(1,1,2)$ vs $(2,1,2)$. En ambos $f_2=T$; $f_1$ cambia de $F$ a $T$.
  \item Mayor $f_2$: $(1,2,1)$ vs $(2,1,2)$. En ambos $f_1=T$; $f_2$ cambia de $F$ a $T$.
\end{itemize}

\paragraph{$P_{10}$: $(\textit{triOut}=3) \wedge (s_2+s_3>s_1)$.}
\begin{itemize}
  \item Mayor $g_1$: $(1,1,2)$ vs $(1,2,2)$. En ambos $g_2=T$; $g_1$ cambia.
  \item Mayor $g_2$: $(2,1,1)$ vs $(1,2,2)$. En ambos $g_1=T$; $g_2$ cambia.
\end{itemize}

\paragraph{Predicados unitarios.} La misma suite produce, en distintos tests, ambos valores de
$P_2$, $P_3$, $P_4$, $P_5$ y $P_7$. Por ejemplo, $P_5$ vale $T$ en $(1,2,3)$ (sin lados iguales)
y $F$ en $(1,1,1)$.

\subsection{Implementacion JUnit}

La logica de instrumentacion vive en \texttt{TriTypTraceSupport}: para cada entrada genera la
secuencia de predicados alcanzados con sus clausulas y, al final, compara la salida
\textit{trackeada} con el valor real devuelto por \texttt{TriTyp.triang}. La suite
\texttt{TriTypExercise4CaccTest} verifica:

\begin{itemize}
  \item las salidas esperadas de \texttt{triang} para los 14 casos;
  \item los pares CACC para $P_1, P_6, P_8, P_9, P_{10}$ usando un \texttt{assertCaccPair} que
  exige cambio de la clausula mayor, igualdad de las menores y cambio del predicado;
  \item que la suite alcance ambos valores de $P_2, P_3, P_4, P_5$ y $P_7$.
\end{itemize}

Ejecucion (desde \texttt{tp5/assignment-5-rodeghiero}):

\begin{lstlisting}[style=codestyle]
JAVA_HOME=$(/usr/libexec/java_home -v 17) mvn -Dmaven.repo.local=.m2 \
    -Djacoco.skip=true test
\end{lstlisting}

con resultado \texttt{BUILD SUCCESS} y todos los tests verdes.

\section*{Conclusiones}
\addcontentsline{toc}{section}{Conclusiones}

A lo largo del practico se aplico la jerarquia de cobertura logica $\text{PC},\text{CC}$,
$\text{CACC},\text{RACC}$ sobre predicados de complejidad creciente:

\begin{itemize}
  \item El \textbf{Ejercicio 1} muestra que CC y PC son criterios \emph{distintos}: no se
  subsumen entre si, y se pueden construir suites que satisfacen uno pero no el otro siempre
  que el predicado admita variabilidad apropiada.
  \item El \textbf{Ejercicio 2} expone la diferencia entre cobertura estructural (aristas) sobre
  un CFG expandido y cobertura logica (CACC) sobre el predicado original: ambos suben en la
  jerarquia desde Edge Coverage, pero por caminos distintos, y un test suite valido para uno
  puede no serlo para el otro.
  \item El \textbf{Ejercicio 3} traslada el analisis a un caso real (\texttt{Thermostat}). La
  diferencia entre CC y PC se ve claramente y la suite CACC obliga a configurar el ambiente
  para que cada clausula \emph{efectivamente} determine la salida.
  \item El \textbf{Ejercicio 4} ilustra que en un metodo con varios predicados encadenados
  (\texttt{TriTyp.triang}) hay que razonar predicado por predicado y reaprovechar casos para
  cumplir CACC en cada uno, manteniendo trazabilidad respecto a los valores reales que devuelve
  el codigo.
\end{itemize}

La instrumentacion y los asserts vinculan el modelo logico con el comportamiento real del
programa, evitando la discordancia entre lo razonado en papel y lo que ocurre en la JVM.

\end{document}
